\section{Number Theory}
\label{sec:numbertheory}

The previous section computed \emph{numbers}: quantities carried to whatever precision
you asked for, where the last digit is always an approximation and the question is only
how good an approximation. This section returns to the opposite pole---the exact
arithmetic of the whole numbers, where nothing rounds and nothing overflows. An integer
in Mathilda is not a 64-bit machine word straining against its limits but, the moment it
outgrows one, a GMP\index{GMP} big integer that grows without bound, so
\(\B{Fibonacci}[200]\) and \(2^{1000}\) are computed to the last digit as readily as
\(2+2\). Number theory is the mathematics of exactly these objects, and it is where a
computer algebra system shows a face very unlike a calculator's: not \emph{how much} but
\emph{how}---how an integer factors, which integers are prime, how they behave under a
modulus, how a rational unfolds into a continued fraction.

We take the subject in the classical order. First the greatest common divisor and the
Euclidean algorithm that computes it (\B{GCD}, \B{LCM}, \B{ExtendedGCD}, \B{CoprimeQ},
\B{Divisible}); then primes and how Mathilda tests, counts, and finds them (\B{PrimeQ},
\B{Prime}, \B{PrimePi}, \B{NextPrime}); then the factorization pipeline that breaks a
composite into primes (\B{FactorInteger}, \B{IntegerExponent}, \B{SquareFreeQ}); the
classical arithmetic functions built on that factorization (\B{Divisors},
\B{DivisorSigma}, \B{EulerPhi}, \B{MoebiusMu}, \B{LiouvilleLambda}, \B{PrimeOmega},
\B{PrimeNu}); modular arithmetic and the structure of the residues (\B{PowerMod},
\B{MultiplicativeOrder}, \B{PrimitiveRoot}, \B{JacobiSymbol}); continued fractions
(\B{ContinuedFraction}, \B{FromContinuedFraction}); the digit and combinatorial functions
(\B{IntegerDigits}, \B{Binomial}, \B{PartitionsP}); and we close by assembling almost all
of it into a working RSA cryptosystem.

\subsection{The Euclidean algorithm and B\'ezout}
\label{ssec:nt-euclid}

The greatest common divisor is the foundation stone. \B{GCD} returns the largest integer
dividing all of its arguments and \B{LCM}\index{least common multiple} the smallest they all divide; both take any
number of arguments and both were met in passing in \S\ref{sec:arithmetic}. What makes
them a starting point rather than a footnote is the algorithm underneath---Euclid's, the
oldest nontrivial algorithm still in daily use\index{Euclidean algorithm}---and the
identity that ties the two functions together:

\begin{codepairs}
\pair{number-theory/euclid/1}{The largest integer dividing both \(24\) and \(36\).}
\pair{number-theory/euclid/2}{And the smallest that \(4\) and \(6\) both divide.}
\pair{number-theory/euclid/3}{Their product, \(\gcd\cdot\operatorname{lcm}\), \dots}
\pair{number-theory/euclid/4}{\dots equals the product of the numbers themselves---a
classic identity, true for any pair.}
\end{codepairs}

\B{ExtendedGCD} does more than report the gcd: it returns a \emph{B\'ezout
certificate}\index{Bezout's identity@B\'ezout's identity}, integers expressing the gcd as
a linear combination of the inputs. For \(\B{ExtendedGCD}[a,b]=\{g,\{s,t\}\}\) you always
have \(sa+tb=g\):

\begin{codepairs}
\pair{number-theory/euclid/5}{Here \(\gcd(240,46)=2\), with the certificate coefficients
\(s=-9\), \(t=47\).}
\pair{number-theory/euclid/6}{Read in the right order they check out: \((-9)\cdot240 +
47\cdot46 = 2\). These coefficients are exactly what inverts one number modulo another---
the engine behind the modular division of \S\ref{ssec:nt-modular}.}
\pair{number-theory/euclid/7}{\B{CoprimeQ}\index{coprime} asks the yes/no question directly: \(14\) and
\(15\) share no factor.}
\pair{number-theory/euclid/8}{\(14\) and \(21\) share a \(7\), so they are not coprime.}
\pair{number-theory/euclid/9}{\B{Divisible} tests one-way divisibility: \(4\mid100\).}
\pair{number-theory/euclid/10}{\(7\) does not divide \(100\).}
\end{codepairs}

\usagebox{ExtendedGCD}

\calloutnote{Under the hood}{Euclid's algorithm replaces the pair \((a,b)\) with
\((b,\,a\bmod b)\) until the second entry is zero; the extended form carries the two
B\'ezout coefficients along the same recurrence. Mathilda does not run this in
interpreted code---it hands the integers to GMP's \texttt{mpz\_gcd} and
\texttt{mpz\_gcdext}, whose subquadratic ``half-GCD'' handles thousand-digit inputs in
microseconds. The partial quotients Euclid discards on the way are, as
\S\ref{ssec:nt-contfrac} will show, precisely a continued fraction~\cite{knuth-taocp2}.}

\subsection{Primes and primality}
\label{ssec:nt-primes}

A prime is an integer greater than \(1\) with no divisors but \(1\) and itself, and the
primes are the multiplicative atoms from which every integer is built. \B{PrimeQ} is the
test:

\begin{codepairs}
\pair{number-theory/primes/1}{\(97\) is prime.}
\pair{number-theory/primes/2}{\(91=7\cdot13\) is not---a number that looks prime until you
try to factor it.}
\pair{number-theory/primes/3}{\B{Prime} indexes the primes directly: the hundredth prime
is \(541\).}
\pair{number-theory/primes/4}{\B{PrimePi} counts them: there are \(25\) primes up to
\(100\).}
\end{codepairs}

The counting function \(\pi(x)\) is where number theory meets analysis. The Prime Number
Theorem\index{Prime Number Theorem} says \(\pi(x)\sim x/\ln x\), and Mathilda lets you
watch the approximation at \(x=10^{6}\):

\begin{codepairs}
\pair{number-theory/primes/5}{The exact count below a million.}
\pair{number-theory/primes/6}{The theorem's estimate \(x/\ln x\)---low by about eight
percent here, an error the theorem promises shrinks (slowly) to zero.}
\pair{number-theory/primes/7}{\B{NextPrime} finds the next prime above its argument,
instantly even past a million.}
\pair{number-theory/primes/8}{A negative second argument walks the other way, so this is
the largest prime \emph{below} \(100\)---there is no separate \texttt{PreviousPrime}.}
\end{codepairs}

\usagebox{PrimeQ}
\usagebox{PrimePi}

\B{PrimePi} is not one algorithm but a family, selectable through \texttt{Method}. Below a
threshold Mathilda simply sieves; above it, counting every prime is hopeless and it uses
a combinatorial method that never enumerates them---all agreeing, of course, on the
answer:

\begin{codepairs}
\pair{number-theory/primes/9}{Legendre's sieve-free recurrence.}
\pair{number-theory/primes/10}{The modern Del\'eglise--Rivat method---same count,
reachable to far larger \(x\).}
\end{codepairs}

A last twist: the notion of ``prime'' depends on the ring. In the Gaussian
integers\index{Gaussian integers} \(\mathbb{Z}[i]\) a rational prime may split, and
\B{PrimeQ} tests this too under \texttt{GaussianIntegers -> True}:

\begin{codepairs}
\pair{number-theory/primes/11}{\(5=(2+i)(2-i)\) is \emph{not} a Gaussian prime---it splits.}
\pair{number-theory/primes/12}{\(3\) stays prime in \(\mathbb{Z}[i]\); it is inert.}
\pair{number-theory/primes/13}{And \(2+3i\) is a Gaussian prime, its norm \(13\) being a
rational prime.}
\end{codepairs}

\calloutnote{Under the hood}{\B{PrimeQ} is not trial division. For machine-sized integers
it runs a deterministic Miller--Rabin test over a fixed witness set; for larger integers
it uses the Baillie--PSW test (a strong Fermat test composed with a Lucas test), for which
no counterexample is known below the ranges Mathilda will ever see. Both are
probabilistic in theory and, at these sizes, certain in practice~\cite{crandallpomerance}.}

\subsection{Integer factorization}
\label{ssec:nt-factor}

Testing primality is easy; \emph{factoring} a composite is the hard problem on which much
of modern cryptography rests. \B{FactorInteger} returns the prime factorization as a list
of \(\{\text{prime},\text{exponent}\}\) pairs:

\begin{codepairs}
\pair{number-theory/factorization/1}{\(360=2^{3}\cdot3^{2}\cdot5\).}
\pair{number-theory/factorization/2}{The fifth Fermat number \(2^{32}+1\). Fermat
conjectured all such numbers prime; Euler found this factor \(641\) in
\(1732\)\index{Fermat number}, and Mathilda recovers his result at once.}
\pair{number-theory/factorization/3}{A pandigital composite peeled apart into eight primes.}
\pair{number-theory/factorization/4}{A \(19\)-digit semiprime---two ten-digit primes---
factored in milliseconds. Its size is trivial for Mathilda but hints at the wall that RSA
hides behind: multiply the primes to hundreds of digits and no known method finishes.}
\end{codepairs}

\usagebox{FactorInteger}

The single call above is a whole cascade of algorithms, and \texttt{Method} lets you
name one explicitly:

\begin{codepairs}
\pair{number-theory/factorization/5}{The same semiprime, factored by Pollard's rho method
on demand\index{Pollard rho method}.}
\pair{number-theory/factorization/6}{\B{IntegerExponent} reads off the multiplicity of a
base directly: \(10^{6}\parallel10^{6}\).}
\pair{number-theory/factorization/7}{The power of \(2\) dividing \(360\).}
\pair{number-theory/factorization/8}{\B{SquareFreeQ}\index{squarefree}: \(30=2\cdot3\cdot5\) has no repeated
prime.}
\pair{number-theory/factorization/9}{\(12=2^{2}\cdot3\) does, so it is not square-free.}
\end{codepairs}

\calloutnote{Under the hood}{\B{FactorInteger} runs a pipeline, each stage suited to a
different size of factor. It strips small primes by trial division, then hunts medium
factors with Pollard's rho (Brent's variant) and the \(p\pm1\) methods, which succeed when
\(p-1\) or \(p+1\) is smooth. For the genuinely hard factors it turns to Lenstra's
elliptic-curve method (ECM)\index{elliptic curve method}, whose running time depends on
the size of the factor found rather than the number
factored~\cite{lenstraecm,crandallpomerance}, with continued-fraction and quadratic-sieve
routines behind that. The default chooses among them by size; \texttt{Method} overrides
the choice.}

\calloutnote{Performance}{The asymmetry between \S\ref{ssec:nt-primes} and this section is
the whole point. Deciding that a \(300\)-digit number is composite takes milliseconds;
producing its factors is, as far as anyone knows, infeasible. RSA
(\S\ref{ssec:nt-rsa}) is exactly this gap, weaponised.}

\subsection{Arithmetic functions}
\label{ssec:nt-arith}

Once an integer is factored, a whole family of \emph{arithmetic
functions}\index{arithmetic function}---functions of a single integer, defined through its
factorization---becomes computable~\cite{hardywright}. Most are \emph{multiplicative}\index{multiplicative
function}: their value on \(mn\) is the product of their values on \(m\) and \(n\) when
\(m,n\) are coprime. The most basic list or count the divisors:

\begin{codepairs}
\pair{number-theory/arithmetic-functions/1}{\B{Divisors} lists every positive divisor of
\(36\), in order.}
\pair{number-theory/arithmetic-functions/2}{\B{DivisorSigma} with first argument \(0\)
counts them: \(36\) has nine divisors.}
\pair{number-theory/arithmetic-functions/3}{With first argument \(1\) it sums them;
\(\sigma_1(28)=56\), a fact we return to below.}
\end{codepairs}

Euler's totient \(\varphi(n)\) counts the integers in \(1,\dots,n\) coprime to \(n\)---
equivalently, the size of the group of units modulo \(n\):

\begin{codepairs}
\pair{number-theory/arithmetic-functions/4}{\(\varphi(36)=12\).}
\pair{number-theory/arithmetic-functions/5}{The same number from the product formula
\(n\prod_{p\mid n}(1-1/p)\)---this is what ``multiplicative'' buys you.}
\end{codepairs}

The remaining functions read off the shape of the factorization. The M\"obius function
\(\mu(n)\)\index{Mobius function@M\"obius function} is \(0\) when a square divides \(n\)
and \(\pm1\) otherwise; Liouville's \(\lambda(n)\), \(\Omega(n)\), and \(\nu(n)\) count
prime factors with and without multiplicity:

\begin{codepairs}
\pair{number-theory/arithmetic-functions/6}{\(\mu(30)=(-1)^{3}\): three distinct primes,
none repeated.}
\pair{number-theory/arithmetic-functions/7}{\(\mu(12)=0\): the square \(4\) divides it.}
\pair{number-theory/arithmetic-functions/8}{\(\lambda(12)=(-1)^{\Omega(12)}=(-1)^{3}=-1\).}
\pair{number-theory/arithmetic-functions/9}{\B{PrimeOmega} counts primes \emph{with}
multiplicity: \(360=2^3\cdot3^2\cdot5\) has \(\Omega=6\).}
\pair{number-theory/arithmetic-functions/10}{\B{PrimeNu} counts \emph{distinct} primes:
\(\nu(360)=3\).}
\pair{number-theory/arithmetic-functions/11}{Like the divisor functions, these extend to
\(\mathbb{Z}[i]\): the Gaussian divisors of \(5\) exhibit its splitting \(5=(2+i)(2-i)\).}
\end{codepairs}

\usagebox{DivisorSigma}
\usagebox{EulerPhi}
\usagebox{MoebiusMu}

These functions are not curiosities; the divisor sum is the gateway to one of the oldest
questions in mathematics. A \emph{perfect number}\index{perfect number} equals the sum of
its proper divisors---equivalently, \(\sigma_1(n)=2n\)---and Mathilda finds the first four
by brute force in an instant:

\begin{codepairs}
\pair{number-theory/perfect/1}{Every perfect number below \(10^{4}\): \(6,28,496,8128\).}
\pair{number-theory/perfect/2}{Euclid's formula \(2^{p-1}(2^{p}-1)\) for \(p=2,3,5,7\)
reproduces exactly that list\index{Euclid--Euler theorem}.}
\end{codepairs}

The match is the Euclid--Euler theorem: \(2^{p-1}(2^{p}-1)\) is perfect precisely when
\(2^{p}-1\) is a \emph{Mersenne prime}\index{Mersenne prime}. So the even perfect numbers
are in exact correspondence with the Mersenne primes, and testing one is testing the
other:

\begin{codepairs}
\pair{number-theory/perfect/3}{\(2^{13}-1=8191\) is a Mersenne prime, giving the fifth
perfect number.}
\pair{number-theory/perfect/4}{But \(2^{11}-1\) is \emph{not} prime---the exponent being
prime is necessary, not sufficient.}
\pair{number-theory/perfect/5}{Indeed \(2^{11}-1=2047=23\cdot89\).}
\end{codepairs}

\subsection{Modular arithmetic and congruences}
\label{ssec:nt-modular}

To work modulo \(n\) is to care only about remainders. \B{Mod} (met in
\S\ref{sec:arithmetic}, and always returning a residue with the sign of the modulus) is
the entry point; the number-theoretic weight sits in what you can do inside a modulus. The
first essential is fast modular exponentiation, because forming a huge power and
\emph{then} reducing it is wasteful when you can reduce at every step:

\begin{codepairs}
\pair{number-theory/modular/1}{Forming \(7^{100}\) in full and reducing---correct, but it
builds an \(85\)-digit integer first.}
\pair{number-theory/modular/2}{\B{PowerMod} gives the same residue without ever forming
the giant power.}
\pair{number-theory/modular/3}{A far larger exponent, still immediate.}
\end{codepairs}

\B{PowerMod} carries two overloads that pack a surprising amount of number theory. A
negative exponent asks for the modular \emph{inverse}, and a fractional one for a modular
\emph{root}:

\begin{codepairs}
\pair{number-theory/modular/4}{The inverse of \(3\) modulo \(7\) is \(5\), since
\(3\cdot5=15\equiv1\).}
\pair{number-theory/modular/5}{A modular square root: \(3^{2}=9\equiv2\pmod 7\), computed
by the Tonelli--Shanks algorithm\index{Tonelli--Shanks algorithm}.}
\end{codepairs}

The \emph{multiplicative order} of \(a\) modulo \(n\) is the least power returning \(a\) to
\(1\); when it reaches the full \(\varphi(n)\), \(a\) is a \emph{primitive root}\index{primitive
root}, a generator of the units:

\begin{codepairs}
\pair{number-theory/modular/6}{\(2\) has order \(3\) modulo \(7\) (\(2^{3}=8\equiv1\)).}
\pair{number-theory/modular/7}{\(3\) has order \(6\)---the full group---so it is a
primitive root.}
\pair{number-theory/modular/8}{\B{PrimitiveRoot} returns the least one, \(3\).}
\pair{number-theory/modular/9}{\B{PrimitiveRootList} gives all of them; there are
\(\varphi(\varphi(7))=\varphi(6)=2\).}
\pair{number-theory/modular/10}{Euler's theorem in one line: \(3^{\varphi(10)}\equiv1\),
because \(3\) is coprime to \(10\).}
\end{codepairs}

The order also frames the hardest question in this circle of ideas. Reversing it---given a
base \(a\) and a target \(b\), find the exponent \(x\) with \(a^{x}\equiv b\pmod n\)---is the
\emph{discrete logarithm problem}\index{discrete logarithm problem}, and the three-argument
form of \B{MultiplicativeOrder} solves exactly that, returning the least exponent whose power
lands in a given residue set:

\begin{codepairs}
\pair{number-theory/modular/13}{The discrete logarithm of \(5\) to base \(3\) modulo \(7\):
the least \(x\) with \(3^{x}\equiv5\), namely \(x=5\) (indeed \(3^{5}=243\equiv5\)).}
\end{codepairs}

\calloutnote{Theory}{Computing the order of \(a\) is easy---it divides \(\varphi(n)\) and is
found by stripping prime factors from it. The discrete logarithm is the \emph{reverse}
question, and it is believed \emph{hard}: no polynomial-time algorithm is known over a
general prime field, and that presumed hardness is the foundation of Diffie--Hellman key
exchange and the ElGamal and DSA signature schemes\index{discrete logarithm problem!in
cryptography}. The best general methods---baby-step giant-step, Pollard's rho for logarithms,
Pohlig--Hellman (which reduces the problem to the prime-power factors of the group order),
and index calculus---are surveyed in~\cite{hac}. Mathilda's three-argument form simply walks
the orbit \(a^{1},a^{2},\dots\), so it is a teaching tool for small moduli, not a break of
the problem at cryptographic sizes.}

Finally, the \B{JacobiSymbol}\index{Jacobi symbol} \(\left(\tfrac{a}{n}\right)\) encodes quadratic
residues\index{quadratic reciprocity}---whether \(a\) is a square modulo \(n\):

\begin{codepairs}
\pair{number-theory/modular/11}{\(\left(\tfrac{2}{7}\right)=+1\): \(2\) is a square mod
\(7\) (indeed \(3^{2}\equiv2\), as above).}
\pair{number-theory/modular/12}{\(\left(\tfrac{3}{7}\right)=-1\): \(3\) is a non-residue.}
\end{codepairs}

\usagebox{PowerMod}
\usagebox{MultiplicativeOrder}
\usagebox{PrimitiveRoot}
\usagebox{JacobiSymbol}

\calloutnote{Theory}{Two theorems govern this section. Fermat's little
theorem\index{Fermat's little theorem}---\(a^{p-1}\equiv1\pmod p\) for a prime \(p\nmid
a\)---is the special case \(n=p\) of \emph{Euler's theorem}\index{Euler's theorem},
\(a^{\varphi(n)}\equiv1\pmod n\) whenever \(\gcd(a,n)=1\). Together with the fact that the
units modulo a prime form a \emph{cyclic} group---generated by a primitive root---they are
the whole basis of what follows, and of RSA in \S\ref{ssec:nt-rsa}~\cite{irelandrosen}.}

Mathilda has no single \texttt{ChineseRemainder} builtin, and it is instructive to see
that none is needed: the Chinese Remainder Theorem\index{Chinese Remainder Theorem} is a
short composition of the tools above. Given residues modulo pairwise-coprime moduli, we
reconstruct the unique solution modulo their product with cofactors and \B{PowerMod}
inverses---the classical recipe of Sun Zi:

\begin{codepairs}
\pair{number-theory/crt/3}{The combined modulus \(3\cdot5\cdot7=105\).}
\pair{number-theory/crt/4}{Each cofactor \(N/n_i\).}
\pair{number-theory/crt/5}{Each cofactor's inverse modulo its own \(n_i\).}
\pair{number-theory/crt/6}{The weighted sum, reduced---the unique residue \(23\).}
\pair{number-theory/crt/7}{And it checks: \(23\) leaves remainders \(2,3,2\) modulo
\(3,5,7\), as required.}
\end{codepairs}

\subsection{Continued fractions}
\label{ssec:nt-contfrac}

A continued fraction\index{continued fraction} writes a number as
\(a_0+\cfrac{1}{a_1+\cfrac{1}{a_2+\cdots}}\), a form that captures rational approximation
better than any decimal. \B{ContinuedFraction} produces the list of partial quotients and
\B{FromContinuedFraction} folds it back:

\begin{codepairs}
\pair{number-theory/continued-fractions/1}{\(415/93=[4;2,6,7]\). These quotients are
exactly Euclid's from \S\ref{ssec:nt-euclid}, run on \(415\) and \(93\).}
\pair{number-theory/continued-fractions/2}{Folding the list back recovers the rational
exactly---the two operations are inverse.}
\end{codepairs}

Irrationals have \emph{infinite} expansions, so you ask for a fixed number of terms---and
the classical patterns appear at once:

\begin{codepairs}
\pair{number-theory/continued-fractions/3}{\(\sqrt2=[1;\overline{2}]\), purely periodic.}
\pair{number-theory/continued-fractions/4}{\(\sqrt7\) has period \(1,1,1,4\).}
\pair{number-theory/continued-fractions/5}{Asked without a length, Mathilda returns the
period in a nested list---every quadratic irrational is \emph{eventually periodic}, which
is Lagrange's theorem\index{quadratic irrational}.}
\pair{number-theory/continued-fractions/6}{\(\pi\) shows no pattern, but its early terms
are the key to good approximations.}
\end{codepairs}

Truncating \(\pi\)'s expansion just after the large term \(15\) and folding it back gives
the approximation known since antiquity:

\begin{codepairs}
\pair{number-theory/continued-fractions/7}{\(355/113\), from only four terms.}
\pair{number-theory/continued-fractions/8}{It agrees with \(\pi\) to better than
\(3\times10^{-7}\)---the best rational approximation with any denominator under a thousand.}
\end{codepairs}

\usagebox{ContinuedFraction}
\usagebox{FromContinuedFraction}

\calloutnote{Theory}{A large partial quotient signals an unusually good rational
approximation just before it: truncating there makes the discarded tail
\(1/(a_{k+1}+\cdots)\) tiny. \(\pi\)'s term \(292\) is why \(355/113\) is so
extraordinarily accurate. The convergents obtained by truncation are, in a precise sense,
the \emph{best} rational approximations to a number\index{continued fraction!convergent}~\cite{khinchin}---the theory that also governs Pell's
equation, which \S\ref{ssec:alg-diophantine} solves through the continued fraction of
\(\sqrt D\).}

\subsection{Digits and bases}
\label{ssec:nt-digits}

The written form of an integer---its digits in a given base---is itself a number-theoretic
object, and Mathilda exposes it directly. \B{IntegerDigits} and \B{FromDigits} convert both
ways, in any base:

\begin{codepairs}
\pair{number-theory/digits/1}{The base-ten digits of \(123456\).}
\pair{number-theory/digits/2}{In base two, \(255\) is eight ones.}
\pair{number-theory/digits/3}{In base sixteen it is two ``fifteens''.}
\pair{number-theory/digits/4}{\B{FromDigits} rebuilds the integer from its base-two digits.}
\pair{number-theory/digits/5}{\B{IntegerString} renders the digits as text, here the
familiar hexadecimal \texttt{ff}.}
\pair{number-theory/digits/6}{\B{IntegerLength} counts the digits without listing them.}
\end{codepairs}

The digit sum is the basis of the schoolroom rule of \emph{casting out nines}---a number
is congruent to its digit sum modulo \(9\), because \(10\equiv1\):

\begin{codepairs}
\pair{number-theory/digits/7}{\B{DigitSum} adds the base-ten digits: \(1+2+3+4+5+6=21\).}
\pair{number-theory/digits/8}{Its residue modulo \(9\) \dots}
\pair{number-theory/digits/9}{\dots equals the number's own residue modulo \(9\). That is
casting out nines.}
\pair{number-theory/digits/10}{\B{DigitCount} tallies how often each digit occurs---one
\(1\), two \(2\)s, three \(3\)s, four \(4\)s (the final slot counts the zeros).}
\end{codepairs}

\usagebox{IntegerDigits}

\subsection{Combinatorial number theory}
\label{ssec:nt-combinatorial}

Number theory shades into combinatorics through the factorials, binomial coefficients, and
partition functions---all computed exactly, all growing explosively. The factorial family
comes first:

\begin{codepairs}
\pair{number-theory/combinatorial/1}{\(20!\), exact to the last of its nineteen digits.}
\pair{number-theory/combinatorial/2}{The double factorial \(9!!=9\cdot7\cdot5\cdot3\cdot1\).}
\pair{number-theory/combinatorial/3}{The falling factorial \(10\cdot9\cdot8\), via
\B{FactorialPower}.}
\pair{number-theory/combinatorial/4}{\B{Binomial}: the number of \(3\)-subsets of a
\(10\)-set.}
\pair{number-theory/combinatorial/5}{Reduced modulo \(7\) it is \(1\)---as Lucas's theorem
predicts from the base-\(7\) digits of \(10\) and \(3\)\index{Lucas's theorem}.}
\end{codepairs}

The partition functions count the ways to write \(n\) as a sum of positive integers.
\B{IntegerPartitions} lists them and \B{PartitionsP} counts them, and the count grows so
fast that its exact evaluation is a small miracle of analysis:

\begin{codepairs}
\pair{number-theory/combinatorial/6}{The seven partitions of \(5\).}
\pair{number-theory/combinatorial/7}{Their count, \(p(10)=42\).}
\pair{number-theory/combinatorial/8}{\(p(100)=190569292\)---already far past listing.}
\pair{number-theory/combinatorial/9}{\(p(1000)\), a \(32\)-digit integer computed exactly.}
\pair{number-theory/combinatorial/10}{\B{PartitionsQ} counts partitions into
\emph{distinct} parts.}
\end{codepairs}

Finally the Fibonacci and Lucas numbers, the archetypal integer sequences:

\begin{codepairs}
\pair{number-theory/combinatorial/11}{\(F_{50}\).}
\pair{number-theory/combinatorial/12}{\(F_{200}\), a \(42\)-digit number, still instant.}
\pair{number-theory/combinatorial/13}{The tenth Lucas number.}
\pair{number-theory/combinatorial/14}{A jewel of an identity: \(\gcd(F_m,F_n)=F_{\gcd(m,n)}\),
here \(\gcd(F_{12},F_{18})=F_{6}\).}
\end{codepairs}

\usagebox{Binomial}
\usagebox{PartitionsP}
\usagebox{Fibonacci}

\calloutnote{Under the hood}{\B{PartitionsP} does not count partitions one by one---that
would be hopeless at \(n=1000\). It evaluates the Hardy--Ramanujan--Rademacher
series\index{Hardy--Ramanujan formula}~\cite{andrews}, a convergent sum of the analytic terms that gives
\(p(n)\) exactly after enough of them, rounded to the nearest integer. \B{Fibonacci} uses
fast doubling---computing \(F_{2k}\) and \(F_{2k+1}\) from \(F_k\) and \(F_{k+1}\)---so
\(F_n\) costs \(O(\log n)\) big-integer multiplications rather than \(n\) additions. The
close relatives \B{BernoulliB}, \B{EulerE}, and \B{Zeta}, whose values encode still deeper
number theory, belong to the special functions of \S\ref{sec:specialfns}.}

\subsection{Putting it together: textbook RSA}
\label{ssec:nt-rsa}

Almost every tool of this section meets in one place: the RSA cryptosystem\index{RSA},
whose security is exactly the factoring gap of \S\ref{ssec:nt-factor}. We use the historic
textbook primes \(p=61\), \(q=53\) so every step is legible; the session is continuous,
each line building on the last. First the public modulus and its secret totient:

\begin{codepairs}
\pair{number-theory/rsa/1}{Pick the first prime, \(p=61\)\dots}
\pair{number-theory/rsa/2}{\dots and the second, \(q=53\).}
\pair{number-theory/rsa/3}{The modulus \(n=pq=3233\) is public; its factors are the secret.}
\pair{number-theory/rsa/4}{The totient \(\varphi(n)=3120\)\dots}
\pair{number-theory/rsa/5}{\dots which for a product of two primes is \((p-1)(q-1)\), as
knowing the factorization lets us confirm.}
\end{codepairs}

Next a public exponent coprime to \(\varphi(n)\), and the private exponent that inverts it:

\begin{codepairs}
\pair{number-theory/rsa/6}{The conventional public exponent \(e=17\)\dots}
\pair{number-theory/rsa/7}{\dots coprime to \(\varphi(n)\), so it is invertible.}
\pair{number-theory/rsa/8}{\B{ExtendedGCD} already exposes the inverse in its B\'ezout
coefficient: \((-367)\cdot17+2\cdot3120=1\).}
\pair{number-theory/rsa/9}{The private key \(d\) is that inverse as a positive residue,
\(2753\) (indeed \(-367+3120\)).}
\pair{number-theory/rsa/10}{Confirming the defining relation \(ed\equiv1\pmod{\varphi(n)}\).}
\end{codepairs}

Now encryption raises the message to \(e\) modulo \(n\), and decryption raises the
ciphertext to \(d\):

\begin{codepairs}
\pair{number-theory/rsa/11}{The plaintext number \(42\).}
\pair{number-theory/rsa/12}{Encrypted, it becomes \(2557\)---bearing no visible relation
to the original.}
\pair{number-theory/rsa/13}{Decryption with the private exponent returns \(42\) exactly.}
\pair{number-theory/rsa/14}{The round trip closes.}
\end{codepairs}

Why it works is Euler's theorem from \S\ref{ssec:nt-modular}: because \(ed\equiv1\pmod{\varphi(n)}\),
raising the message to \(e\) and then to \(d\) raises it to a power \(\equiv1\), which
leaves it unchanged modulo \(n\). Anyone holding \((n,e)\) can encrypt; only the holder of
\(d\)---recoverable only by factoring \(n\)---can decrypt. Every step ran on integers no
larger than a few thousand, but nothing about the method changes when the primes have a
hundred digits each; only the factoring stays out of reach.

\subsection*{Where this connects}

Number theory is the exact spine of Mathilda's arithmetic. Its foundation is the integer
model of \S\ref{sec:arithmetic}---the same floored \B{Mod} and GMP big integers, here put
to work. Its natural sequel is the arithmetic of \emph{solutions}: the classical Diophantine
questions---Pell's equation, Pythagorean triples, sums of squares, and the rest---are
answered by \B{Solve} and \B{Reduce} over the integers, treated in
\S\ref{ssec:alg-diophantine}, where the continued fractions of \S\ref{ssec:nt-contfrac}
reappear as the engine behind Pell. The polynomial \B{Factor} of \S\ref{sec:algebra} is the
algebraic mirror of \B{FactorInteger}, square-free decomposition and Hensel lifting
standing in for trial division and rho. And the integer-valued special functions that shade
number theory into analysis---\B{BernoulliB}, \B{Zeta}, and their kin---await in
\S\ref{sec:specialfns}. Underneath every result here is GMP, and the same exact integers
feed the compiler of Chapter~\ref{ch:compilation} when a number-theoretic loop needs to run
at machine speed.
